\section{Numerical evaluation of properties}
\label{sec:state}
\subsection{Association as an implicit calculation}
% P:state-01
A density-specified state takes temperature, molar density and normalized, charge-neutral species composition.
For an associating fluid, the native calculation first solves the mass-action system of Eq.~\eqref{eq:mass-action}.
Its unbonded-site fractions change with composition and molecular parameters, so their implicit response must enter both properties and fitting columns.

% P:state-02
The native association iteration starts at unit site fractions, with an analytic two-population estimate for sufficiently strong association.
It solves the analytic mass-action Jacobian by full-pivot LU decomposition and damps a Newton step to preserve $0<X_A\le1$.
When the Jacobian does not separate its pivots sufficiently, substitution supplies the update.
The mass-action residual, rather than the step length alone, decides convergence.
Table~\ref{tab:defaults} gives the default residual and work limits.
The implicit derivative calculation subsequently uses the unweighted defining Jacobian at the converged fractions; a singular or ill-conditioned system makes the affected derivative unavailable.
No composition floor is introduced to make that derivative appear regular.

\subsection{Pressure closure and mechanical branches}
% P:state-03
A pressure-specified state adds the scalar equation
\begin{equation}
P(T,\rho,\mathbf{x})=\rho RT\left[1+\rho\left(\frac{\partial\ares}{\partial\rho}\right)_{T,\mathbf{x}}\right]=P_{\mathrm{target}}.
\label{eq:pressure-closure}
\end{equation}
Pressure, its slope and, when required for root location, its curvature follow from the same energy~\citep{Gross2001}.
The density search operates in a bounded packing domain and first inspects the liquid side.
Resolved pressure brackets with positive slope can then identify a liquid root and a lower-density vapor root.
The requested branch selects the highest or lowest resolved stable root; an unspecified branch is ambiguous when several stable roots remain.
An anchor from a neighboring state instead defines a local window in which the root must be unique.
These choices matter at subcritical conditions because the same temperature, pressure and composition can have several density roots.

% P:state-04
Slope and curvature samples can shorten root location when they are consistent with a single mechanical loop.
Unresolved shape information returns the search to a finite logarithmic density scan, with spinodal bisection separating the stable pressure brackets.
When a sample contradicts that assumed shape, the interval is subdivided, and a root is accepted only if successive finer subdivisions show the same sign changes of $P-P_{\mathrm{target}}$.
Safeguarded Newton refinement remains inside each pressure bracket.
With $r=(P-P_{\mathrm{target}})/(\rho RT)$ and bracket width $\Delta\rho$, termination requires $|r|\le\tau_\rho$ and $\Delta\rho\,|\mathrm{d}r/\mathrm{d}\rho|\le\tau_\rho$, with a finite nonzero derivative; an exactly zero computed residual is also accepted.
An unresolved bracket at the refinement budget is reported as a failure.
A small reduced residual does not imply the same relative pressure error at low compressibility factor: the latter is the reduced residual divided by $P_{\mathrm{target}}/(\rho RT)$.
Association failures interrupt the relevant scan and preserve the failed density and cause, so a disconnected set of evaluable points cannot silently become a complete root search.

\begin{algorithm}[tbp]
\caption{Evaluation of a density- or pressure-specified state}
\label{alg:state}
\begin{algorithmic}[1]
\Require Selected parameters, $T$, $\mathbf{x}$, density or pressure, branch and requested quantities
\State Admit units, support, charge, topology and required thermal intervals.
\State Select the contributions and the derivative recording (Section~\ref{sec:derivatives}) for these species and this temperature interval.
\State At each trial density solve Eq.~\eqref{eq:mass-action}, preserving $0<X_A\le1$.
\State Stop the site solve only when its residual meets the requested criterion; otherwise report its cause.
\If{pressure is specified}
\State Locate liquid-side evidence and the requested pressure bracket using analytic slopes and curvature.
\State Use the logarithmic scan and spinodal subdivision when structured location is unresolved.
\State Refine stable crossings by safeguarded Newton; require both reduced-residual and bracket-width criteria, or an exactly zero computed residual.
\State Select the declared extreme root, or the unique root in the supplied anchor window.
\EndIf
\State Evaluate the requested ideal, residual or total quantities at this one state.
\State Differentiate association and pressure closure implicitly for requested actions.
\State Return values, achieved residuals and conditioning; keep failed quantities explicitly unavailable.
\end{algorithmic}
\end{algorithm}

% P:state-05
Algorithm~\ref{alg:state} returns a mechanically admitted state, distinguished from a globally stable phase.
Supporting Information Fig.~\ref{fig:si-F2-density-roots-1} shows the roots and spinodals found for propane and aqueous NaCl and the pressure closure of the benchmark requests.
Gibbs-energy comparisons across compositions belong to Section~\ref{sec:equilibrium}; finite density sampling cannot exclude every arbitrarily narrow additional loop.


\subsection{Property assembly and reuse}
% P:state-06
Property requests reuse the converged density, sites and required energy derivatives.
Selected quantities compute only their dependencies; total calorics consume thermal records, while residual quantities remain available without them.
A new support or thermal interval selects a corresponding recording of the same equations and refreshes the numerical state.
